\chapter{Introduction}

\begin{comment}

    Describe problem(s) that we are trying to solve.
     - complex (changing) dynamics of aerospace vehicles.
     - fault tolerant control.
     - accurate model needed
     - mav
       • model not possible due to
         · complex nature of mav dynamics
         · inconsistent manufacturing
         · parameter adaptation during time
         · susceptible to unforseen disturbances.        
     - new technological improvements spurred the development of innovative more advanced aerospace systems.
       • From 

    
    Why it's relevant to the aerospace industry.
     - aerospace vehicle, duh, maybe this should go with the first item

    What the basic solution to this problem could be, current solutions
     - (adaptive) non-linear control
    
    Research question

    Proposed solution of the preliminary thesis and why.
     - RL
     - Short description of what RL is capable of
     - Give examples of previous research or real world application (keep short or out)

    Short description on the approach that will be taken during the thesis to use reinforcement learning to solve the problem.
     - Policy gradient
     - Self tuning quad-copter
     - Simulation

    Describe the contents in this report and how it will be used during the actual thesis.
     - Introduction into reinforcement learning
     - Tabular methods
     - Approximate solutions 
     - History of RL. Past, Present and Future (Or maybe just the latest in RL)
     - Thesis plan of attack


\end{comment}


\todo{Rewrite this paragraph so it's easier to read.}
% From the rise of Unmanned Aerial Vehicles to increasingly congested air airspaces, 
The technological advancements and steadily increasing demand for air and space travel of the past decades have lead to the development of new, more complex aerospace systems, resulting in new challenges in the field of aerospace control. At the same time advancements in automated control have been seen in many domains, ranging from industrial manufacturing and household appliances to the rapid widespread adoption of UAVs. However most of these systems are only able to complete their tasks within the narrow, predictable conditions they where designed for. Autonomous control systems need to be able to adapt to unexpected situations and environments. This has motivated researchers to search for new innovative methods that can be used to solve these new problems. 

Adaptive control methods allow autonomous systems to reconfigure them self in case of changing conditions. This opens the possibility to create fault tolerant control systems capable of safely recovering after a failure

One of the leading causes of fatal air traffic accidents is loss of control. If fully autonomous UAV's are to be used in the future, advanced flight control systems capable of adapting to rapidly changing conditions are necessary. This is because survival can only only be guaranteed if sta




